\documentclass[12pt,a4paper]{amsart}

\usepackage{amsmath,amssymb,amsthm}
\usepackage{mathtools}
\usepackage[utf8]{inputenc}
\usepackage[ngerman]{babel}
\usepackage{hyperref}
\usepackage{geometry}
\geometry{margin=2.5cm}

% -------------------------------------------------------
% Theorem-Umgebungen
% -------------------------------------------------------
\newtheorem{theorem}{Satz}[section]
\newtheorem{lemma}[theorem]{Lemma}
\newtheorem{corollary}[theorem]{Korollar}
\newtheorem{proposition}[theorem]{Proposition}
\theoremstyle{definition}
\newtheorem{definition}[theorem]{Definition}
\newtheorem{remark}[theorem]{Bemerkung}

% -------------------------------------------------------
% Eigene Befehle
% -------------------------------------------------------
\newcommand{\N}{\mathbb{N}}
\newcommand{\Z}{\mathbb{Z}}
\newcommand{\R}{\mathbb{R}}
\newcommand{\C}{\mathbb{C}}
\DeclareMathOperator{\Re}{Re}
\DeclareMathOperator{\Li}{Li}

% -------------------------------------------------------
\begin{document}
% -------------------------------------------------------

\title{Riemanns explizite Formel und der Primzahlsatz}

\author{Michael Fuhrmann}
\address{Specialist Mathematics Project, Build 84}
\email{specialist-maths@localhost}
\date{\today}

\begin{abstract}
Wir entwickeln Riemanns explizite Formel, die die Primzahlz\"ahlfunktion
$\pi(x)$ mit den Nullstellen von $\zeta(s)$ verbindet:
\[
  \pi(x) = \Li(x) - \sum_\rho \Li(x^\rho) - \log 2
  + \int_x^{\infty} \frac{dt}{t(t^2-1)\log t},
\]
wobei $\Li(x) = \int_2^x dt/\log t$ und die Summe \"uber alle nicht-trivialen
Nullstellen $\rho$ l\"auft.
Aus dieser Formel leiten wir den Primzahlsatz
$\pi(x) \sim x/\log x$ ab und -- unter Annahme der Riemann-Hypothese --
den scharfen Fehlerterm $|\pi(x) - \Li(x)| < \tfrac{1}{8\pi}\sqrt{x}\,\log x$
f\"ur $x \ge 2657$.
Die von-Mangoldt-Funktion $\Lambda(n)$ und die Tschebyschew-Funktionen
$\psi(x)$ und $\theta(x)$ werden eingef\"uhrt und mit $\zeta$ in Beziehung gesetzt.
Abschlie\ss end weisen wir auf die Verbindung zu Dirichlet-$L$-Funktionen
und der Verallgemeinerten Riemann-Hypothese (GRH) hin:
Unter GRH gelten analoge Fehlerterme f\"ur Primzahlen in arithmetischen
Progressionen, und die L-Nullstellen bestimmen die Gleichverteilung.
\end{abstract}

\maketitle
\tableofcontents

% ================================================================
\section{Einleitung}
% ================================================================

Die Br\"ucke zwischen analytischen Eigenschaften von $\zeta(s)$ und
arithmetischen Eigenschaften der Primzahlen ist die \emph{explizite Formel}.
Riemann skizzierte sie 1859 \cite{Riemann1859};
sie wurde von von~Mangoldt \cite{vonMangoldt1895} und Hadamard \cite{Hadamard1896}
rigoros begr\"undet.

Die Grundidee: Die Lage der Nullstellen von $\zeta(s)$ kodiert pr\"azise
Information \"uber die Primzahlverteilung.
Die Riemann-Hypothese ($\Re(\rho) = 1/2$) liefert den sch\"arfstm\"oglichen
Fehlerterm.

% ================================================================
\section{Arithmetische Funktionen und Tschebyschew-Funktionen}
% ================================================================

\begin{definition}[Von-Mangoldt-Funktion]
\label{def:mangoldt}
Die \emph{von-Mangoldt-Funktion} ist definiert durch
\[
  \Lambda(n) \;=\;
  \begin{cases}
    \log p & \text{falls } n = p^k \text{ f\"ur eine Primzahl } p
             \text{ und } k \ge 1, \\
    0      & \text{sonst.}
  \end{cases}
\]
\end{definition}

\begin{definition}[Tschebyschew-Funktionen]
\label{def:chebyshev}
Die \emph{erste Tschebyschew-Funktion} ist
$\theta(x) = \sum_{p \le x} \log p$,
und die \emph{zweite Tschebyschew-Funktion} (\emph{Psi-Funktion}) ist
\[
  \psi(x) = \sum_{n \le x} \Lambda(n) = \sum_{p^k \le x} \log p
           = \sum_{p \le x} \lfloor\log_p x\rfloor \cdot \log p.
\]
\end{definition}

\begin{proposition}[Beziehungen zwischen $\pi$, $\theta$, $\psi$]
\label{prop:psi_theta_pi}
\begin{align}
  \psi(x) &= \theta(x) + \theta(x^{1/2}) + \theta(x^{1/3}) + \cdots
  \label{eq:psi_theta} \\
  \pi(x) &= \frac{\theta(x)}{\log x}
    + \int_2^x \frac{\theta(t)}{t(\log t)^2}\,dt.
  \label{eq:pi_theta}
\end{align}
Insbesondere sind $\pi(x) \sim x/\log x$, $\psi(x) \sim x$ und
$\theta(x) \sim x$ \"aquivalent.
\end{proposition}

\begin{proof}
F\"ur~\eqref{eq:psi_theta}: Jeder Term $\theta(x^{1/k})$ z\"ahlt
$\sum_{p \le x^{1/k}} \log p$, also die Beitr\"age der Primzahlpotenzen
$p^k \le x$.
F\"ur~\eqref{eq:pi_theta}: Abel'sche Summation
$\pi(x) = \theta(x)/\log x + \int_2^x \theta(t)/[t(\log t)^2]\,dt$,
was durch partielle Integration aus $d[\pi(t)] = d[\theta(t)/\log t]$ folgt.
\end{proof}

% ================================================================
\section{Logarithmische Ableitung und \texorpdfstring{$\psi$}{psi}}
% ================================================================

\begin{theorem}[Logarithmische Ableitung von $\zeta$]
\label{thm:log_deriv}
F\"ur $\Re(s) > 1$ gilt
\[
  -\frac{\zeta'(s)}{\zeta(s)} \;=\; \sum_{n=1}^{\infty} \frac{\Lambda(n)}{n^s}.
\]
\end{theorem}

\begin{proof}
Logarithmieren des Euler-Produkts $\zeta(s) = \prod_p (1-p^{-s})^{-1}$:
\[
  \log\zeta(s) = -\sum_p \log(1 - p^{-s}) = \sum_p \sum_{k=1}^{\infty} \frac{p^{-ks}}{k}.
\]
Gliedweises Differenzieren (f\"ur $\Re(s)>1$ gerechtfertigt):
\[
  \frac{\zeta'(s)}{\zeta(s)} = -\sum_p \sum_{k=1}^{\infty} \frac{\log p}{p^{ks}}
  = -\sum_{n=1}^{\infty} \frac{\Lambda(n)}{n^s}. \qedhere
\]
\end{proof}

% ================================================================
\section{Perrons Formel und die explizite Formel f\"ur \texorpdfstring{$\psi$}{psi}}
% ================================================================

\begin{theorem}[Perrons Formel, vereinfacht]
\label{thm:perron}
F\"ur $c > 1$ und $x > 0$ keine Primzahlpotenz:
\[
  \psi(x) = \frac{1}{2\pi i} \int_{c-i\infty}^{c+i\infty}
  \left(-\frac{\zeta'(s)}{\zeta(s)}\right) \frac{x^s}{s}\,ds.
\]
\end{theorem}

\begin{proof}[Beweisskizze]
Schl\"usselidentit\"at: F\"ur $c > 0$ und $x > 0$:
\[
  \frac{1}{2\pi i} \int_{c-i\infty}^{c+i\infty} \frac{x^s}{s}\,ds
  = \begin{cases} 1 & x > 1, \\ 0 & 0 < x < 1. \end{cases}
\]
Einsetzen der Dirichlet-Reihe f\"ur $-\zeta'/\zeta$ und gliedweises Integrieren
(durch Konvergenzabsch\"atzungen gerechtfertigt) ergibt die Behauptung.
Vollst\"andiger Beweis in \cite{Davenport2000}, Kapitel~17.
\end{proof}

\begin{theorem}[Explizite Formel f\"ur $\psi(x)$]
\label{thm:explicit_psi}
F\"ur $x > 1$ und $x$ keine Primzahlpotenz gilt:
\begin{equation}
  \label{eq:explicit_psi}
  \psi(x) \;=\; x - \sum_\rho \frac{x^\rho}{\rho}
  - \frac{\zeta'(0)}{\zeta(0)} - \tfrac{1}{2}\log(1 - x^{-2}),
\end{equation}
wobei die Summe \"uber alle nicht-trivialen Nullstellen $\rho$ von $\zeta$
l\"auft (symmetrisch nach $|\Im(\rho)| \le T$ geordnet, dann $T \to \infty$;
bedingt konvergent in dieser Anordnung).

Es gilt $\zeta'(0)/\zeta(0) = \log(2\pi) \approx 1{,}8379\ldots$\,.
\end{theorem}

\begin{proof}[Beweisskizze]
Man verschiebt die Integrationskontur in Perrons Formel nach links und
sammelt die Residuen der Polstellen und Nullstellen von $-\zeta'(s)/\zeta(s)$:
\begin{itemize}
  \item Residuum bei $s = 1$: $\zeta$ hat einen einfachen Pol,
    also hat $-\zeta'/\zeta$ einen einfachen Pol mit Residuum $-1$;
    Beitrag: $x$.
  \item Residuen bei nicht-trivialen Nullstellen $\rho$: Beitrag $-x^\rho/\rho$.
  \item Triviale Nullstellen $s = -2m$ ($m=1,2,3,\ldots$):
    $-\zeta'/\zeta$ hat bei jedem $s = -2m$ einen einfachen Pol mit Residuum $+1$.
    Perrons Formel liefert den Beitrag $+x^{-2m}/(-2m)$ f\"ur jedes $m$.
    Summation ergibt:
    $\displaystyle\sum_{m=1}^\infty \frac{x^{-2m}}{-2m}
     = -\tfrac{1}{2}\sum_{m=1}^\infty \frac{(x^{-2})^m}{m}
     = \tfrac{1}{2}\log(1-x^{-2})$
    (unter Verwendung von $\log(1-u)=-\sum u^m/m$).
    In der Endformel erscheint dieser Term als $-\tfrac{1}{2}\log(1-x^{-2})$;
    da $\log(1-x^{-2}) < 0$ f\"ur $x > 1$, ist $-\tfrac{1}{2}\log(1-x^{-2}) > 0$
    eine kleine positive Korrektur.
  \item Polstelle bei $s = 0$: Beitrag $-\zeta'(0)/\zeta(0)$.
\end{itemize}
Vollst\"andiger Beweis in \cite{Davenport2000}, Kapitel~17.
\end{proof}

% ================================================================
\section{Der Primzahlsatz}
% ================================================================

\begin{theorem}[Primzahlsatz]
\label{thm:pnt}
\[
  \pi(x) \;\sim\; \frac{x}{\log x} \quad (x \to \infty).
\]
\"Aquivalent dazu: $\psi(x) \sim x$.
\end{theorem}

\begin{proof}[Beweisskizze via expliziter Formel]
\textbf{Schritt 1} (Keine Nullstellen auf $\Re(s) = 1$):
Mertens (1898) beobachtete: F\"ur $\sigma > 1$ und reelles $t$:
\[
  \log|\zeta^3(\sigma)\,\zeta^4(\sigma+it)\,\zeta(\sigma+2it)| \ge 0.
\]
Dies folgt aus $3 + 4\cos\theta + \cos 2\theta = 2(1+\cos\theta)^2 \ge 0$.
W\"are $\zeta(1+it_0) = 0$, so g\"alte f\"ur $\sigma \to 1^+$:
$\log|\zeta^3(\sigma)| \approx +3\log\tfrac{1}{\sigma-1}$ (einfacher Pol),
aber $\log|\zeta^4(\sigma+it_0)| \le -4\log\tfrac{1}{\sigma-1} + O(1)$
(Nullstelle mindestens erster Ordnung).
Da $4 > 3$, g\"inge die Gesamtsumme gegen $-\infty$ -- Widerspruch zu $\ge 0$.
Also $\zeta(1+it) \ne 0$ f\"ur $t \ne 0$.

\textbf{Schritt 2} (PZS aus nullstellenfreiem Gebiet):
Sobald $\zeta(s) \ne 0$ auf $\Re(s) = 1$ bekannt ist, zeigt die explizite
Formel~\eqref{eq:explicit_psi}, dass der dominante Term $x$ ist
und alle anderen Terme $o(x)$ sind.
Also $\psi(x) \sim x$ und damit $\pi(x) \sim x/\log x$.
\end{proof}

\begin{remark}
Dieser Ansatz ist eine vereinfachte Version der urspr\"unglichen Beweise
von Hadamard \cite{Hadamard1896} und de~la~Vall\'ee~Poussin \cite{Poussin1896},
die den Primzahlsatz unabh\"angig voneinander 1896 bewiesen.
Der elementare Beweis (ohne komplexe Analysis) stammt von
Erd\H{o}s und Selberg (1949).
\end{remark}

% ================================================================
\section{Fehlerterme: Die Rolle der Nullstellen}
% ================================================================

\begin{theorem}[Fehlerterm in Abh\"angigkeit von den Nullstellen]
\label{thm:error_term}
Unter der Annahme, dass alle nicht-trivialen Nullstellen von $\zeta$
Realteil $\le \beta_0 < 1$ haben, gilt
\[
  \psi(x) = x + O\!\bigl(x^{\beta_0}(\log x)^2\bigr).
\]
Insbesondere gibt die Riemann-Hypothese ($\beta_0 = \tfrac{1}{2}$):
\[
  \psi(x) = x + O(\sqrt{x}\,\log^2 x).
\]
\end{theorem}

\begin{proof}[Beweisskizze]
In der expliziten Formel~\eqref{eq:explicit_psi} tr\"agt die Nullstellensumme
$\sum_\rho x^\rho/\rho$ h\"ochstens $x^{\beta_0}\sum_\rho |\rho|^{-1}$ bei.
Da die Anzahl der Nullstellen mit $|\rho| \le T$ von der Gr\"o\ss enordnung
$O(T\log T)$ ist, konvergiert $\sum_\rho |\rho|^{-1}$ mit logarithmischer
Divergenz, und die gesamte Summe ist $O((\log T)^2)$ durch sorgf\"altige
Absch\"atzung.
\end{proof}

\begin{theorem}[Schoenfelds explizite Schranke unter RH]
\label{thm:schoenfeld}
Unter der Annahme der Riemann-Hypothese gilt f\"ur alle $x \ge 2657$:
\[
  |\pi(x) - \Li(x)| \;<\; \frac{1}{8\pi}\,\sqrt{x}\,\log x.
\]
\end{theorem}

\begin{proof}[Beweisskizze]
Schoenfeld \cite{Schoenfeld1976} bewies diese Schranke durch sorgf\"altige
Auswertung der Nullstellensumme $\sum_{|\gamma| \le T} x^{1/2+i\gamma}/(1/2+i\gamma)$
mit expliziten Konstanten unter Verwendung der RH.
\end{proof}

% ================================================================
\section{Riemanns urspr\"ungliche explizite Formel f\"ur \texorpdfstring{$\pi(x)$}{pi(x)}}
% ================================================================

\begin{theorem}[Riemanns explizite Formel f\"ur $\pi(x)$]
\label{thm:riemann_explicit}
Sei $\Li(x) = \int_2^x dt/\log t$. Dann gilt f\"ur $x > 1$:
\begin{equation}
  \label{eq:riemann_explicit}
  \pi(x) \;=\; \Li(x) - \sum_\rho \Li(x^\rho) - \log 2
  + \int_x^{\infty} \frac{dt}{t(t^2-1)\log t}.
\end{equation}
Die Summe l\"auft \"uber alle nicht-trivialen Nullstellen $\rho$ (geordnet
nach $|\Im(\rho)|$); das Integral konvergiert f\"ur $x > 1$.
\end{theorem}

\begin{proof}[Beweisskizze]
Riemann verwendete die \glqq Primzahlpotenz-Z\"ahlfunktion\grqq
\[
  J(x) = \sum_{n=1}^{\infty} \frac{1}{n}\,\pi(x^{1/n})
  = \frac{1}{2\pi i}\int_{c-i\infty}^{c+i\infty} \frac{\log\zeta(s)}{s}\,x^s\,ds,
\]
aus der $\pi(x)$ durch M\"obius-Inversion gewonnen wird:
$\pi(x) = \sum_{n=1}^\infty \mu(n)/n \cdot J(x^{1/n})$.
Die Pole und Nullstellen von $\log\zeta(s)$ (via Hadamard-Produkt)
liefern dann~\eqref{eq:riemann_explicit}.
Vollst\"andige Herleitung in \cite{Titchmarsh1986}, Kapitel~3.
\end{proof}

% ================================================================
\section{Verbindung zu \texorpdfstring{$L$}{L}-Funktionen und GRH}
% ================================================================

\begin{definition}[Dirichlet-L-Funktionen und Primzahlen in Progressionen]
\label{def:pnt_progressions}
F\"ur $\gcd(a,q) = 1$ sei
\[
  \pi(x; q, a) = \#\{p \le x : p \equiv a \pmod q\}.
\]
Durch die Orthogonalit\"at der Dirichlet-Charaktere gilt:
\[
  \pi(x; q, a) = \frac{1}{\phi(q)}\sum_{\chi \bmod q} \bar\chi(a)\,
  \pi(x, \chi),
\]
wobei $\pi(x,\chi)$ die Nullstellen von $L(s,\chi)$ enth\"alt.
\end{definition}

\begin{theorem}[Dirichlet'scher Satz mit Fehlerterm unter GRH]
\label{thm:dirichlet_grh}
Unter Annahme der GRH f\"ur alle $L(s,\chi)$ mit $\chi \pmod q$ gilt:
\[
  \pi(x; q, a) = \frac{\Li(x)}{\phi(q)}
  + O\!\Bigl(\sqrt{x}\,\log(qx)\Bigr).
\]
Insbesondere sind Primzahlen gleichm\"a\ss ig auf die $\phi(q)$
Restklassen prim zu $q$ verteilt.
\end{theorem}

% ================================================================
\section{Schlussfolgerung}
% ================================================================

Riemanns explizite Formel ist die zentrale Verbindung zwischen den
Nullstellen von $\zeta(s)$ und der Primzahlverteilung.
Der Primzahlsatz folgt daraus, dass $\zeta$ keine Nullstellen auf
$\Re(s)=1$ hat.
Unter der Riemann-Hypothese ist der Fehlerterm $O(\sqrt{x}\log x)$ --
im Wesentlichen das sch\"arfstm\"ogliche Resultat.
Die Formel~\eqref{eq:riemann_explicit} zeigt: Jede nicht-triviale Nullstelle
$\rho$ leistet einen Beitrag $\Li(x^\rho) \approx x^\rho/\log x$ zu $\pi(x)$;
je mehr Nullstellen abseits der kritischen Geraden l\"agen, desto gr\"o\ss er
w\"aren diese Schwingungen.

% ================================================================
\begin{thebibliography}{10}

\bibitem{Riemann1859}
B.~Riemann,
\emph{\"Uber die Anzahl der Primzahlen unter einer gegebenen Gr\"o\ss e},
Monatsberichte der Berliner Akademie (1859), 671--680.

\bibitem{vonMangoldt1895}
H.~von~Mangoldt,
Zu Riemann's Abhandlung ``\"Uber die Anzahl der Primzahlen \ldots'',
\emph{J.\ Reine Angew.\ Math.}\ \textbf{114} (1895), 255--305.

\bibitem{Hadamard1896}
J.~Hadamard,
Sur la distribution des z\'eros de la fonction $\zeta(s)$ et ses
cons\'equences arithm\'etiques,
\emph{Bull.\ Soc.\ Math.\ France}\ \textbf{24} (1896), 199--220.

\bibitem{Poussin1896}
C.-J.\ de~la~Vall\'ee~Poussin,
Recherches analytiques sur la th\'eorie des nombres premiers,
\emph{Ann.\ Soc.\ Sci.\ Bruxelles}\ \textbf{20} (1896), 183--256.

\bibitem{Davenport2000}
H.~Davenport,
\emph{Multiplicative Number Theory}, 3.~Aufl.,
Springer, New York, 2000.

\bibitem{Titchmarsh1986}
E.~C.~Titchmarsh,
\emph{The Theory of the Riemann Zeta-Function}, 2.~Aufl.,
Oxford University Press, 1986.

\bibitem{Schoenfeld1976}
L.~Schoenfeld,
Sharper bounds for the Chebyshev functions $\theta(x)$ and $\psi(x)$, II,
\emph{Math.\ Comp.}\ \textbf{30} (1976), 337--360.

\end{thebibliography}

\end{document}
